\documentclass[12pt]{article}
\usepackage[utf8]{inputenc}
\usepackage[T1]{fontenc}
\usepackage{hyperref}
\title{Northern Pame Language Resources for K-12 Curriculum Development}
\author{Compiled by Claude Sonnet 5.5 (claude-sonnet-5-5) for Andr\'as Kornai}
\date{}
\begin{document}
\maketitle

\section{Introduction}
Northern Pame (Xi'iuy; ISO 639-3 pmq; Glottocode nort2984) is an Oto-Pamean language of San Luis Potosí, Mexico. Speaker counts vary between sources (about 6,000 on Wikipedia, far fewer in ASJP). Academic documentation exists; school materials were not found in search.

\section{Primary Educational Resources}
No primers or textbooks were found in search results.

\section{Linguistic and Reference Materials}
\begin{itemize}
\item Scott Berthiaume, \emph{A Phonological Grammar of Northern Pame} (2003 dissertation; later SIL ebook), bibliographic entry: \url{https://phoible.org/sources/pmq_berthiaume2003}
\item \emph{Diccionario xi'iuy (pame norte) de los municipios de Tamasopo y Rayón, San Luis Potosí} (2007, IIH; only cited in a search summary, no link to the dictionary itself was seen).
\item Overview of Pame languages, Wikipedia: \url{https://en.wikipedia.org/wiki/Pame_languages}
\item Pame chapter (PDF), Colegio de México CELL: \url{https://cell.colmex.mx/archivos/QWxpbQogNDAKZG9jdW1lbnRv/39%20Pame.pdf}
\item ASJP word list entry: \url{https://asjp.clld.org/languages/NORTHERN_PAME}
\end{itemize}

\section{Religious and Community Resources}
None found.

\section{Digital and Audio Resources}
\begin{itemize}
\item The Northern Pame Collection of Clifton Pye, Scott Berthiaume and Barbara Pfeiler, AILLA (recordings and transcriptions of child speech): \url{https://islandora-ailla.lib.utexas.edu/islandora/object/ailla%3A268837}
\end{itemize}

\section{Curriculum Development Context}
Existing material is mainly phonological and archival. A curriculum could build on the dictionary and the AILLA recordings, but reading materials and teacher resources would need to be created.
\end{document}